\paragraph{Maximal runs and ranks.}\label{ta:maximal-ranks}
Let $R=\Reach_\sigma(Q_0)$ be computed without expanding goals. A maximal run is infinite, or finite and ending at a goal or a state with no policy successor. Assume $R$ is finite, all its states are safe, goals are terminal, and every non-goal has a successor. If the reachable non-goal subgraph has a cycle, first follow a path from a root to that cycle and then repeat it. Every edge is policy-permitted, including its selected nondeterministic outcome, so this is an infinite maximal run avoiding goals. Universal finite completion excludes it without invoking fairness.

Conversely, if that subgraph is acyclic, no path can visit a non-goal twice. Every path therefore has at most $|R\setminus\Goal|$ non-goal edges before ending. A non-goal endpoint is impossible by the successor premise. Thus every maximal run ends at a goal. Define $\rank(q)$ as the maximum length of a goal-ending path from $q$, and set goal ranks to zero. Such a path exists because the graph is finite, acyclic off goals, and has no non-goal sinks. For each edge $q\to q'$ from a non-goal, prefixing a longest path from $q'$ proves $\rank(q)\geq1+\rank(q')$. This gives a positive natural-number rank at non-goals. Finally, a strict decrease around a cycle would give $\rank(q)>\rank(q)$, a contradiction. These implications establish all three equivalences. Without the successor premise, a non-goal sink is an acyclic counterexample.

